\section{Resultados}
\label{sec:resultados}

Los resultados se presentan en cinco subsecciones (A--E) más un estudio de calibración del directorio.
Todos los experimentos utilizan el mismo banco de 80 consultas descrito en la Sección~\ref{sec:banco} y los parámetros fijados en la Sección~\ref{sec:parametros}, salvo que se indique lo contrario.

% ---------------------------------------------------------------
\subsection{A. Protocolo completo}
\label{sec:proto-completo}

El banco de 80 consultas se procesó siguiendo el pipeline completo descrito en la Sección~\ref{sec:metodologia}: llenado con 200 imágenes por clase, formación del directorio con interacciones intercaladas, y evaluación sobre el banco semántico.

\paragraph{Fase temprana.}
El sistema alcanzó una \textbf{precisión temprana del 97.5\%} con una tasa de rechazo del 1.25\%.
El rechazo no viene de un filtro de vocabulario: cada token con representación fastText entra como pista, y es la propia memoria la que rechaza cuando ningún agente obtiene soporte. El único rechazo del banco corresponde a una consulta sin pista representable, no a una palabra ``fuera de lista''.

\paragraph{Fase madura con lectura corregida (B1).}
Con la lectura de directorio dividida por el conteo de registros de cada agente (B1), la \textbf{precisión madura fue del 98.75\%}.
Los conteos del directorio transactivo al final del experimento fueron $[78, 65, 53]$ para los agentes de \texttt{apple}, \texttt{horse} y \texttt{car}, respectivamente, con una entropía de \textbf{1.567 bits} sobre un máximo teórico de 1.585 bits ($\log_2 3$), lo que corresponde a un balance del 98.9\%.

\paragraph{Fase madura con lectura cruda.}
Sin corrección por conteo, la precisión madura cayó al \textbf{53.75\%}, una diferencia de \textbf{45 puntos porcentuales} respecto a la lectura B1.
Esta diferencia muestra que la especialización real produce masas acumuladas desiguales; sin normalizar, el agente que más interacciones registró domina la puntuación independientemente del contenido semántico de la consulta.

\paragraph{Fidelidad temprana-madura.}
La concordancia entre la clase ganadora en fase temprana y en fase madura fue del \textbf{97.5\%}, lo que indica que ambas fases se apoyan mutuamente: cuando la fase temprana es correcta, la fase madura lo confirma; y cuando la fase temprana falla, la madura raramente lo corrige.

\paragraph{Réplica de prueba (10 consultas).}
La réplica con las 10 consultas de depuración originales produjo conteos $[6, 7, 6]$ con \textbf{1 rechazo honesto}: la consulta ``sweet edible thing'' no produce una pista con soporte en ningún agente, así que la memoria la rechaza correctamente sin penalizar la precisión.

\begin{figure}[H]
    \centering
    \includegraphics[width=0.7\textwidth]{distribucion_ganadores.png}
    \caption{Distribución de agentes ganadores en el banco de 80 consultas ($N=80$).
    Los conteos finales del directorio transactivo son $[78, 65, 53]$ para \texttt{apple}, \texttt{horse} y \texttt{car}, con entropía 1.567 bits (98.9\% del máximo).}
    \label{fig:distribucion-ganadores}
\end{figure}

% ---------------------------------------------------------------
\subsection[B. Parámetros nativos iota x kappa]{B. Parámetros nativos $\iota \times \kappa$}
\label{sec:param-iota-kappa}

Se realizó un barrido de cuadrícula con $\iota \in \{0, 0.25, 0.5, 1.0\}$ y $\kappa \in \{0, 0.5, 1.0, 1.5\}$ sobre el banco completo de 80 consultas.
Los resultados se muestran en la Tabla~\ref{tab:grid-iota-kappa}.

\begin{table}[H]
    \centering
    \caption{Precisión y tasa de rechazo del banco de 80 consultas en función de $\iota$ y $\kappa$.
    Las columnas de $\kappa$ son idénticas para cada fila, por lo que se muestra una sola columna representativa.}
    \label{tab:grid-iota-kappa}
    \begin{tabular}{cccc}
        \toprule
        $\iota$ & $\kappa$ (cualquier valor) & Precisión (\%) & Rechazo (\%) \\
        \midrule
        0.00  & 0 -- 1.5 & 97.5 &  1.2 \\
        0.25  & 0 -- 1.5 & 77.5 & 21.2 \\
        0.50  & 0 -- 1.5 & 53.8 & 45.0 \\
        1.00  & 0 -- 1.5 &  0.0 & 100.0 \\
        \bottomrule
    \end{tabular}
\end{table}

\paragraph{Efecto de $\iota$.}
El parámetro de inhibición $\iota$ poda activamente el conocimiento escaso antes de que pueda actuar como señal discriminativa.
Con patrones binarios de 300 bits, las relaciones semánticas son muy ralas, de modo que incluso un valor pequeño ($\iota = 0.25$) elimina conexiones legítimas y reduce la precisión 20 puntos, elevando el rechazo al 21.2\%.
Con $\iota = 1.0$ todas las consultas son rechazadas.

\paragraph{Efecto de $\kappa$.}
El parámetro de umbral $\kappa$ no cambió los resultados en el rango probado: las filas de la Tabla~\ref{tab:grid-iota-kappa} son idénticas para todos los valores de $\kappa$.
El análisis del tensor de activaciones revela que las proyecciones superan la media del tensor en aproximadamente dos órdenes de magnitud, lo que hace que el umbral relativo definido por $\kappa$ nunca sea vinculante.

\paragraph{Redundancia de la calibración por división entre la media.}
Un diagnóstico preliminar de 11 consultas de prueba mostró que la compuerta de containment solo (sin dividir por la media de la memoria) logra un 100\% de precisión, mientras que la combinación compuerta + división por media alcanza un 90.9\%.
Esto confirma que con el llenado por instancias las masas quedan igualadas por construcción y la corrección adicional resulta redundante e incluso contraproducente.

% ---------------------------------------------------------------
\subsection{C. Formación del directorio transactivo}
\label{sec:formacion}

Para estudiar cómo el orden de las interacciones afecta la velocidad de aprendizaje del directorio, se evaluó la precisión madura completa en un \textit{snapshot} tras cada interacción $k$, usando cuatro condiciones de presentación.
El umbral de convergencia se fijó en el 90\%.

\begin{figure}[H]
    \centering
    \includegraphics[width=\textwidth]{formacion_directorio.png}
    \caption{Precisión madura en función del número de interacciones $k$ para las cuatro condiciones de formación del directorio.
    La línea discontinua indica el umbral del 90\%.}
    \label{fig:formacion-directorio}
\end{figure}

\paragraph{Orden intercalado.}
Presentando las imágenes en orden intercalado entre clases, el directorio cruza el umbral del 90\% en \textbf{$k = 13$ interacciones}.

\paragraph{Orden barajado (5 semillas).}
Con presentación aleatoria barajada, la transición ocurre en el rango $k \in [15, 24]$ con una media de \textbf{$k \approx 20$} interacciones.

\paragraph{Orden bloqueado por dominio.}
Cuando las interacciones se presentan agrupadas por clase (todas las de una clase antes de pasar a la siguiente), el umbral no se alcanza hasta \textbf{$k = 67$ interacciones}, ya entrado el tercer bloque.
Esto refleja que el directorio no puede enrutar conocimiento que no ha presenciado todavía.

\paragraph{Control con puntuación cruda.}
Con la puntuación sin corregir, la transición se cruza por primera vez en $k = 14$ pero solo se sostiene desde $k = 66$, y la precisión final es del \textbf{93.8\%} con conteos $[84, 67, 45]$.
Comparado con los conteos del sistema corregido $[78, 65, 53]$, el sesgo de la puntuación cruda deja una estructura de directorio diferente aunque la lectura B1 lo compense en el agregado final.

\paragraph{Convergencia al mismo estado final.}
Todas las condiciones con puntuación corregida convergen al mismo estado final (98.75\% de precisión, conteos $[78, 65, 53]$): el orden de las interacciones afecta la velocidad de aprendizaje del directorio, pero no su destino.

\begin{figure}[H]
    \centering
    \includegraphics[width=0.7\textwidth]{entropia_directorio.png}
    \caption{Entropía del directorio transactivo en función de $k$ para las cuatro condiciones de formación.
    La entropía máxima teórica para tres agentes es $\log_2(3) \approx 1.585$ bits.}
    \label{fig:entropia-directorio}
\end{figure}

% ---------------------------------------------------------------
\subsection{D. Capacidad de llenado}
\label{sec:capacidad}

Se varió el número de imágenes registradas por clase $N \in \{25, 50, 100, 200, 328\}$ utilizando llenado en orden imagen-por-imagen (\textit{imagen-etiqueta}).
La aceptación se determinó por containment inverso con $\xi = 0$.
Los \textit{falsos} son imágenes de otras clases presentadas al agente.

\begin{table}[H]
    \centering
    \caption{Métricas de capacidad en función del número de imágenes registradas $N$ por clase.
    ``Enrutamiento'' es la precisión del directorio de labels; ``Variedad L1'' es la distancia media $L_1$ entre recuperaciones distintas.}
    \label{tab:capacidad}
    \begin{tabular}{ccccc}
        \toprule
        $N$ & Generalización test (\%) & Falsos aceptados (\%) & Enrutamiento labels (\%) & Variedad $L_1$ \\
        \midrule
         25  &  0.0  & 0.0 & 60.0 & 0.00 \\
         50  &  0.0  & 0.0 & 93.8 & 2.25 \\
        100  & 11.7  & 0.0 & 98.8 & 2.53 \\
        200  & 56.7  & 0.0 & 97.5 & 2.75 \\
        328  & 78.3  & 0.0 & 85.0 & 3.09 \\
        \bottomrule
    \end{tabular}
\end{table}

\paragraph{Generalización por cobertura.}
Con $\xi = 0$ el sistema no generaliza desde cero; la generalización aparece a medida que el llenado cubre más del espacio latente de la clase.
A $N = 200$ la tasa de generalización alcanza el 56.7\%, y a $N = 328$ el 78.3\%, sin tener que ajustar $\xi > 0$.

\paragraph{Especificidad perfecta.}
La tasa de falsos aceptados es del \textbf{0\% en todo el rango} de $N$ probado, lo que confirma que el containment inverso discrimina de forma efectiva imágenes fuera de la clase registrada.

\paragraph{Saturación en el enrutamiento de labels.}
A $N = 328$ el enrutamiento de labels cae del 97.5\% (en $N = 200$) al 85.0\%.
Con más latentes por label, las asociaciones cubren una fracción mayor del eje contrario y la compuerta de containment pierde poder discriminativo.
La \textbf{zona de operación recomendada} es $N = 200$.

\begin{figure}[H]
    \centering
    \includegraphics[width=\textwidth]{curva_capacidad.png}
    \caption{Generalización en test, tasa de falsos y enrutamiento de labels en función del número de imágenes registradas $N$.}
    \label{fig:curva-capacidad}
\end{figure}

\begin{figure}[H]
    \centering
    \includegraphics[width=\textwidth]{recall_por_n.png}
    \caption{Cinco recuperaciones decodificadas para la consulta ``fruit'' para cada valor de $N$.
    Con $N$ pequeño las recuperaciones son indistinguibles; la variedad aumenta con la cobertura del llenado.}
    \label{fig:recall-por-n}
\end{figure}

\paragraph{Contraste de llenados: promedio vs.\ instancias.}
Al comparar el llenado por prototipo promedio (un solo vector por clase repetido $N$ veces) contra el llenado por instancias reales, la diferencia es drástica.
Con llenado promedio, 12 muestras de recuperación producen \textbf{1 patrón único} (variedad $L_1 = 0.00$): la memoria colapsa a un atractor único.
Con llenado por instancias, las mismas 12 muestras producen \textbf{12/12 patrones distintos}: el muestreo entrópico opera sobre una distribución real de representaciones y puede explorar el espacio aprendido.
Esta comparación se muestra en la Figura~\ref{fig:variedad-recall}.

\begin{figure}[H]
    \centering
    \includegraphics[width=\textwidth]{variedad_recall.png}
    \caption{Comparación de recuperaciones decodificadas con llenado promedio (izquierda) vs.\ llenado por instancias (derecha).
    Con llenado promedio las recuperaciones son idénticas; con instancias reales aparecen 12 patrones distintos.}
    \label{fig:variedad-recall}
\end{figure}

% ---------------------------------------------------------------
\subsection{E. Hemisferio visual}
\label{sec:hemisferio-visual}

El hemisferio visual opera en dos fases.
La primera es formación, con registro de imágenes del conjunto \texttt{train[200:]}.
La segunda es evaluación, con enrutamiento de imágenes de test mediante el directorio visual.

\paragraph{Fase A: formación con imágenes de \texttt{train[200:]}.}
Se presentaron 384 imágenes (128 por clase) en orden intercalado.
Sobre las imágenes aceptadas, la precisión del agente visual fue del \textbf{100\%}.
La tasa de rechazo fue del \textbf{49.7\%}: con $\xi = 0$ solo se aceptan imágenes suficientemente similares a las ya registradas.
Los conteos del directorio visual al final de la fase A fueron $[91, 71, 31]$ con entropía de \textbf{1.466 bits}.

\paragraph{Fase B: enrutamiento de imágenes de test.}
Las 246 imágenes de test ($82 \times 3$ clases) se enrutaron mediante el directorio visual con lectura B1.
La precisión total fue de \textbf{81/246 = 32.9\%} (imágenes correctamente asignadas sobre el total presentado).
Sobre las imágenes efectivamente aceptadas, la precisión fue del \textbf{100\%}: el directorio visual nunca se equivoca en su asignación.
El rechazo en fase B fue del \textbf{67.1\%}: la cobertura limitada del directorio visual con $\xi = 0$ es la principal restricción.

\paragraph{Evocación de labels: imagen $\to$ palabras.}
Para validar la dirección inversa del hemisferio visual (imagen $\to$ palabras), se calculó la recuperación semántica usando los pesos de la memoria homo-asociativa del latente combinados con la selección top-3 por similitud coseno en el espacio fastText.
De 17 imágenes de evaluación, \textbf{16/17 = 94.1\%} produjeron al menos una etiqueta perteneciente al vocabulario de la clase real.

\paragraph{El rechazo y la salida semántica los decide la memoria.}
Reconocimiento, rechazo y reconstrucción son operaciones de la propia memoria: el encoder ResNet18 solo convierte la imagen en una pista latente, y a partir de ahí la EAM decide. La Figura~\ref{fig:rechazo-imagenes} ilustra el rechazo por containment ($\xi = 0$): imágenes reales que la memoria contiene se aceptan y enrutan, mientras que imágenes reales no contenidas y ejemplos sintéticos fuera de dominio (ruido, color sólido, píxeles barajados) reciben puntuación cero en los tres agentes y se rechazan, sin que intervenga ningún filtro externo. La Figura~\ref{fig:imagen-etiquetas} muestra la salida semántica: para cada imagen aceptada, la reconstrucción es lo que la memoria \emph{evoca} (\texttt{decode} de la recuperación de $M^R_{\mathrm{dom}}$), no la imagen de entrada, y el sistema devuelve las etiquetas asociadas. Que la reconstrucción difiera de la entrada es esperado y honesto: refleja que decide la distribución acumulada en la memoria.

\begin{figure}[H]
    \centering
    \includegraphics[width=0.85\textwidth]{rechazo_imagenes.png}
    \caption{Rechazo de imágenes por la memoria (containment $\xi = 0$). Arriba: imágenes reales aceptadas y enrutadas. En medio: imágenes reales no contenidas, rechazadas. Abajo: ejemplos sintéticos fuera de dominio, también rechazados. El rechazo emerge de la ausencia de soporte, no de un filtro de vocabulario.}
    \label{fig:rechazo-imagenes}
\end{figure}

\begin{figure}[H]
    \centering
    \includegraphics[width=0.85\textwidth]{imagen_a_etiquetas.png}
    \caption{Salida semántica imagen $\to$ etiquetas. Por cada imagen aceptada: entrada (pista), reconstrucción evocada por la memoria (no la entrada) y las tres etiquetas recuperadas, con marca de coincidencia con el vocabulario de la clase real.}
    \label{fig:imagen-etiquetas}
\end{figure}

% ---------------------------------------------------------------
\subsection{Calibración del directorio: estudio de condiciones}
\label{sec:calibracion}

Para determinar la lectura óptima del directorio transactivo se realizó un estudio comparativo de nueve condiciones experimentales con $N = 80$ imágenes registradas, evaluadas sobre el mismo directorio hetero-asociativo del sistema final ($M_{\mathrm{dir}}$, no una aproximación).
Los resultados se muestran en la Tabla~\ref{tab:calibracion}.

\begin{table}[H]
    \centering
    \caption{Precisión madura bajo distintas condiciones de lectura del directorio ($N = 80$).}
    \label{tab:calibracion}
    \begin{tabular}{llc}
        \toprule
        Condición & Descripción & Precisión (\%) \\
        \midrule
        A (cruda)     & Score de activación sin normalizar               & 53.75 \\
        B1            & División entre conteo de registros por agente     & \textbf{98.75} \\
        B2            & División entre raíz cuadrada del conteo           & 75.00 \\
        C             & Cap de número máximo de registros por agente      & 53.75 \\
        D             & Consultas balanceadas entre clases                & 52.82 \\
        E32           & Cuantización binaria en 32 niveles                & 53.75 \\
        E64           & Cuantización binaria en 64 niveles                & 53.75 \\
        F             & Vocabulario curado (sin ruido Apple Inc.)         & 80.00 \\
        G (D+B1+F)    & Consultas balanceadas + B1 + vocabulario curado   & 98.72 \\
        \bottomrule
    \end{tabular}
\end{table}

\begin{figure}[H]
    \centering
    \includegraphics[width=\textwidth]{ablation_heatmap.png}
    \caption{Mapa de calor de precisión madura por condición experimental y número de imágenes $N$.}
    \label{fig:ablation-heatmap}
\end{figure}

\begin{figure}[H]
    \centering
    \includegraphics[width=\textwidth]{ablation_madura.png}
    \caption{Precisión madura en función de $N$ para cada condición de lectura del directorio.}
    \label{fig:ablation-madura}
\end{figure}

\paragraph{Análisis.}
La condición B1 es la que corrige el mecanismo causal: divide la activación acumulada entre el número de registros, obteniendo la \textit{evidencia por interacción}.
Las demás condiciones manipulan síntomas (el vocabulario con F, el sesgo de consultas con D, el número de registros con C) sin atacar la raíz.
La condición G combina B1 con las correcciones de D y F y alcanza el 98.72\%, confirmando que el vocabulario limpio y las consultas balanceadas son beneficiosas pero subordinadas a la corrección principal de B1.

% ---------------------------------------------------------------
\subsection{Costo computacional}
\label{sec:costo}

Se midió el tiempo de pared (\textit{wall-clock}) de reproducir el experimento completo desde cero en una laptop modesta: \textbf{Intel Core i5-13420H} (8 núcleos físicos, 12 hilos), 16~GB de RAM, \textbf{sin uso de GPU} (PyTorch en modo CPU; la tarjeta NVIDIA presente no participa).
La Tabla~\ref{tab:costo} resume los tiempos.

\begin{table}[H]
    \centering
    \caption{Costo computacional de la reproducción completa (CPU, sin GPU). El entrenamiento del encoder es la única pieza costosa; todo el sistema de memorias asociativas opera en segundos.}
    \label{tab:costo}
    \small
    \begin{tabular}{lc}
        \toprule
        Componente & Tiempo \\
        \midrule
        Encoder ResNet18 (50 épocas, 984 imágenes)   & 24.3 min \\
        Sistema asociativo (llenado, fases, directorios, visual) & $<$ 2.5 min \\
        Pipeline completo (etapas 1--8)               & \textbf{26.6 min} \\
        \midrule
        Batería de análisis (estudios A--E + calibración) & $\approx$ 30 min \\
        \bottomrule
    \end{tabular}
\end{table}

El \textbf{entrenamiento del encoder visual} (la única parte que usa gradientes y retropropagación) concentra cerca del 91\% del costo del pipeline.
En contraste, \textbf{todo el sistema de memoria transactiva} —el llenado de las memorias, la fase temprana, la formación de los directorios, el hemisferio visual y la fase madura— suma menos de dos minutos y medio, porque son operaciones matriciales sobre vectores cuantizados, sin entrenamiento.
La batería de análisis (cientos de corridas de evaluación) es trabajo de medición, no parte de un despliegue del sistema.
La lectura es clara: una vez disponible la percepción visual, la maquinaria asociativa que da soporte a la memoria transactiva tiene un costo despreciable y no requiere hardware especializado.
